%%AMC:latex_engine=lualatex
% ↑ AMC 本体にこのファイルを LuaLaTeX で処理させる指定（消さないでください）
%
% 基礎数学 第10回 ベクトルの基礎 確認問題（CC BY-SA 4.0）
\documentclass[a4paper,11pt]{ltjsarticle}
\usepackage{lfs-fonts}
\usepackage{amsmath,amssymb}
\usepackage[nowatermark]{automultiplechoice}
\AMCrandomseed{261001}
\def\AMCbeginQuestion#1#2{\par\noindent{\bfseries 問#1}\quad #2}

\begin{document}

\element{bm10}{
  \begin{question}{BM10-Q1}
    $(1,\, 2) + (3,\, -1)$ は？
    \begin{choiceshoriz}
      \correctchoice{$(4,\, 1)$}
      \wrongchoice{$(4,\, 3)$}
      \wrongchoice{$(-2,\, 3)$}
      \wrongchoice{$(3,\, -2)$}
    \end{choiceshoriz}
  \end{question}
}

\element{bm10}{
  \begin{question}{BM10-Q2}
    $\vec{a} = (3,\, -1)$，$\vec{b} = (1,\, 2)$ のとき $2\vec{a} - \vec{b}$ は？
    \begin{choiceshoriz}
      \correctchoice{$(5,\, -4)$}
      \wrongchoice{$(7,\, 0)$}
      \wrongchoice{$(5,\, 0)$}
      \wrongchoice{$(2,\, -3)$}
    \end{choiceshoriz}
  \end{question}
}

\element{bm10}{
  \begin{question}{BM10-Q3}
    $\vec{a} = (5,\, 12)$ の大きさ $|\vec{a}|$ は？
    \begin{choiceshoriz}
      \correctchoice{$13$}
      \wrongchoice{$17$}
      \wrongchoice{$7$}
      \wrongchoice{$169$}
    \end{choiceshoriz}
  \end{question}
}

\element{bm10}{
  \begin{question}{BM10-Q4}
    $\overrightarrow{\mathrm{AB}} + \overrightarrow{\mathrm{BC}}$ に等しいものは？
    \begin{choiceshoriz}
      \correctchoice{$\overrightarrow{\mathrm{AC}}$}
      \wrongchoice{$\overrightarrow{\mathrm{CA}}$}
      \wrongchoice{$\overrightarrow{\mathrm{BA}}$}
      \wrongchoice{$\vec{0}$}
    \end{choiceshoriz}
  \end{question}
}

\element{bm10}{
  \begin{question}{BM10-Q5}
    2点 A$(2,\, 0)$，B$(8,\, 6)$ を結ぶ線分の中点の座標は？
    \begin{choiceshoriz}
      \correctchoice{$(5,\, 3)$}
      \wrongchoice{$(3,\, 3)$}
      \wrongchoice{$(6,\, 6)$}
      \wrongchoice{$(10,\, 6)$}
    \end{choiceshoriz}
  \end{question}
}

\onecopy{1}{
  \noindent{\large\bfseries 基礎数学　第10回　確認問題}\hfill
  \namefield{氏名\,\underline{\hspace{45mm}}}\par\medskip
  \AMCcodeGridInt[h]{student}{8}
  \medskip
  \insertgroup{bm10}
}

\end{document}
